\chapter{Effects}\label{ch:effects}

Effects are listed in \opt{effects()} separated by blanks; an effect with an argument is written with the argument in parentheses, for example \opt{effects(density contagion gwactivity(0.5) ego(age\_c))}. Below, $a_i\in\{0,1\}$ is the outcome of actor $i$, $x_{ij}$ the tie from $i$ to $j$ (symmetric in an undirected network), $d_i=\sum_j x_{ij}$ the degree and $n_i=\sum_j x_{ij}a_j$ the number of neighbours with the outcome. Every statistic is a sum over the actors with the outcome ($a_i=1$) unless it is described as a count of ties or triples. The names in the first column are those used in \opt{effects()}; the coefficient table prefixes covariate effects with the effect (for example \texttt{ego\_age\_c}). Appendix~\ref{app:translation} gives the functions of \alaamee{} that they correspond to.

\section{Undirected networks}\label{sec:eff-undirected}

\begin{longtable}{@{}p{0.27\textwidth}p{0.70\textwidth}@{}}
\toprule Effect & Statistic \\ \midrule \endhead
\opt{density} & The number of actors with the outcome, $\sum_i a_i$ (the intercept). \\
\opt{activity} & The sum of the degrees of the actors with the outcome, $\sum_i a_i d_i$. \\
\opt{twostar}, \opt{threestar} & $\sum_i a_i\binom{d_i}{2}$ and $\sum_i a_i\binom{d_i}{3}$. \\
\opt{gwactivity}[(\emph{a})] & Geometrically weighted activity, $\sum_i a_i\,e^{-a\,d_i}$ with \emph{a} $>0$; default $\log 2$. \\
\opt{contagion} & The number of ties between two actors with the outcome, $\sum_{i<j}x_{ij}a_ia_j$. \\
\opt{alter2star1} & The partner-activity two-path: $\sum_i a_i\sum_j x_{ij}(d_j-1)$. \\
\opt{alter2star2} & The indirect partner attribute: the number of two-paths between two actors with the outcome. \\
\opt{partneractivity} & The ties between two actors with the outcome, weighted by the degree sum minus two. \\
\opt{partnerpartner} & Two times the number of two-paths with the outcome at all three actors. \\
\opt{triangle1}, \opt{triangle2}, \opt{triangle3} & Triangles weighted by the number of their actors with the outcome (T1), by the number of pairs of them (T2), and only those with all three (T3). \\
\opt{gwcontagion}[(\emph{a})] & $\sum_i a_i\,e^{-a\,n_i}$; default $\log 2$. \\
\opt{logcontagion} & $\sum_i a_i\log(n_i+1)$. \\
\opt{powercontagion}[(\emph{b})] & $\sum_i a_i\,n_i^{1/b}$; default $b=2$. \\
\opt{thresholdcontagion}[(\emph{k})] & The number of actors with the outcome that have at least $k$ neighbours with the outcome, $\sum_i a_i\,[n_i\ge k]$; $k$ a positive integer, default 2. \\
\opt{fractioncontagion}[(\emph{q})] & The number of actors with the outcome of which at least the share $q$ of the neighbours have the outcome, $\sum_i a_i\,[d_i>0,\ n_i/d_i\ge q]$; $0<q\le 1$, default 0.5. \\
\opt{sharecontagion} & $\sum_i a_i\,n_i/d_i$ (0 for an actor without neighbours). \\
\opt{ego(var)} & The covariate of the actors with the outcome, $\sum_i a_i v_i$ (centre it). \\
\opt{alter(var)} & The sum over the actors with the outcome of the covariate of their neighbours. \\
\opt{same(var)}, \opt{diff(var)} & The number of ties of the actors with the outcome to neighbours of the same, or of a different, category. \\
\opt{samecontagion(var)} & The ties between two actors with the outcome and the same category. \\
\opt{samepartneractivity(var)}, \opt{diffpartneractivity(var)} & The two-paths from an actor with the outcome to an actor of the same, or of a different, category. \\
\opt{sameindirect(var)}, \opt{diffindirect(var)} & The two-paths between two actors with the outcome of the same, or of different, categories. \\
\opt{alterbinarytwostar1(var)} & $\sum_i a_i$ times the sum of $(d_j-1)$ over the neighbours with the binary attribute 1. \\
\opt{alterbinarytwostar2(var)} & The two-paths between two actors with the outcome whose middle actor has the binary attribute 1. \\
\bottomrule
\end{longtable}

\paragraph{Exposure thresholds.} \opt{contagion} counts ties between actors with the outcome, so each neighbour with the outcome adds the same amount to the log odds of an actor. The effects \opt{thresholdcontagion}, \opt{fractioncontagion} and \opt{sharecontagion} let the exposure enter in another form: as a step (at least $k$ neighbours, or at least the share $q$ of the neighbours, with the outcome) or as a share. They are the terms with which a specification can represent complex contagion. The listing fits the threshold effect $k=2$ and the share on the karate club, in place of \opt{contagion}.
\showlog{thresholds_karate}
A step function can be almost collinear with \opt{density}: when nearly all actors with the outcome meet the threshold (here 17 of the 18 members of the faction have at least two neighbours with the outcome), the statistic hardly differs from the number of actors with the outcome, and the standard errors are large. Choose $k$ or $q$ so that a good part of the actors with the outcome lies on each side; with a threshold that nearly all or nearly none of them meet, the fit stops with a message that the simulated statistics are collinear (return code 430).

Three more effects use a second source of information and are described in Section~\ref{sec:setting}: \opt{settinghomophily(network)}, \opt{geographichomophily(matname)} and \opt{contagiondist(matname)}.

\section{Directed networks}\label{sec:eff-directed}

In a directed network $x_{ij}$ is the tie from $i$ to $j$. The definitions are those of \alaamee{}.

\begin{longtable}{@{}p{0.27\textwidth}p{0.70\textwidth}@{}}
\toprule Effect & Statistic \\ \midrule \endhead
\opt{density} & The number of actors with the outcome. \\
\opt{sender}, \opt{receiver} & The out-degree and the in-degree of the actors with the outcome. \\
\opt{reciprocity} & The mutual ties of the actors with the outcome. \\
\opt{outstar2}, \opt{outstar3}, \opt{instar2}, \opt{instar3} & Out-stars and in-stars at the actors with the outcome. \\
\opt{mixedstar}, \opt{mixedsource}, \opt{mixedsink} & Mixed two-stars (one in, one out) at the actors with the outcome, and the same with the other end a source or a sink. \\
\opt{contagion} & The ties, in and out, to actors with the outcome; a mutual tie counts twice. \\
\opt{contagionrecip} & The mutual ties between two actors with the outcome. \\
\opt{gwcontagion}[(\emph{a})], \opt{logcontagion}, \opt{powercontagion}[(\emph{b})] & As in undirected networks, with the number of ties (in and out) to actors with the outcome. Their change statistics sum over the out-neighbours with the outcome (counting their in-neighbours with the outcome) and over the in-neighbours with the outcome (counting their out-neighbours with the outcome), as in \alaamee{}. \\
\opt{thresholdcontagion}[(\emph{k})], \opt{fractioncontagion}[(\emph{q})], \opt{sharecontagion} & As in undirected networks; $n_i$ is the number of ties (in and out; a mutual tie counts twice) to actors with the outcome and $d_i$ the number of all ties (in and out) of the actor. \\
\opt{transitive1}, \opt{transitive3} & Transitive triples ($a\to b$, $b\to c$, $a\to c$): the number in which the actor with the outcome is the middle actor $b$; those with all three actors with the outcome. \\
\opt{transitived1}, \opt{transitiveu1} & Transitive triples in which the actor with the outcome is the source $a$, or the sink $c$. \\
\opt{cyclic1}, \opt{cyclic3} & Cyclic triples ($a\to b\to c\to a$) with an actor with the outcome in them; those with all three actors with the outcome. \\
\opt{alterinstar2}, \opt{alteroutstar2} & Two-stars between an actor with the outcome and another actor with the outcome through a common target, or source. \\
\opt{gwsender}[(\emph{a})], \opt{gwreceiver}[(\emph{a})] & Geometrically weighted out- and in-degree. \\
\opt{ego(var)}, \opt{alter(var)} & The covariate of the actor itself, and of its neighbours (in and out). \\
\opt{same(var)}, \opt{diff(var)}, \opt{samecontagion(var)} & Ties (in and out) to actors of the same, or of a different, category; ties to actors with the outcome and the same category. \\
\opt{sendermatch(var)}, \opt{sendermismatch(var)} & Out-ties to actors of the same, or of a different, category. \\
\opt{receivermatch(var)}, \opt{receivermismatch(var)} & In-ties from actors of the same, or of a different, category. \\
\opt{reciprocitymatch(var)}, \opt{reciprocitymismatch(var)} & Mutual ties with actors of the same, or of a different, category. \\
\bottomrule
\end{longtable}

\opt{settinghomophily(network)} and \opt{geographichomophily(matname)} are available as in undirected networks (Section~\ref{sec:setting}).

\section{Two-mode networks}\label{sec:eff-twomode}

The outcome is defined on the actors of one level of a two-mode network, the other level is the set of events, groups or affiliations the actors are tied to (Section~\ref{sec:twomode}). Contagion effects do not exist because there are no ties within a level. The effects are \opt{density}, \opt{activity}, \opt{twostar}, \opt{threestar}, \opt{gwactivity}, \opt{alter2star1}, \opt{alter2star2}, \opt{fourcycle1}, \opt{fourcycle2}, \opt{ego()}, \opt{alter()}, \opt{samepartneractivity()}, \opt{diffpartneractivity()}, \opt{sameindirect()}, \opt{diffindirect()}, \opt{alterbinarytwostar1()}, \opt{alterbinarytwostar2()}, \opt{settinghomophily()} and \opt{geographichomophily()}. The two that exist only here count four-cycles: \opt{fourcycle1} the four-cycles through the actor and another actor of the same level, whatever the outcome of the other actor, and \opt{fourcycle2} those in which the other actor also has the outcome.

\section{Covariates}\label{sec:covariates}

Covariates are variables of the dataset. Continuous covariates are used as they are: centre them (subtract the mean) so that the intercept keeps its meaning, as in the examples. Categorical effects (\opt{same()}, \opt{diff()}, the partner and indirect effects, the match and mismatch effects) compare the values of the variable; a binary attribute (\opt{alterbinarytwostar1()}) must be 0 or 1. A missing covariate value makes the actor count as 0 in the statistics, as a missing outcome does.
